% \chapter{Estratégias de RVD/B}
% \label{sec_cap2_estrategias_rvdb}


% Neste capítulo descreve-se as estratégias empregadas para o processo de RVD/B entre o veículo perseguidor e o alvo.
% O objetivo é traçar o roteiro das abordagens a serem implementadas em uma missão de RVD/B, desde a fase de lançamento até  a manobra orbital de aproximação de proximidade, quando então se dá o acoplamento da nave tipo manipulador robótico na porta de docking da espaçonave alvo, completando assim a operação de RVD/B.
% O Capítulo aborda estratégia para a modelagem matemática e análise dinâmica (Capítulo 2) em conformidade com os conceitos de movimento absoluto e relativo que permeiam as fases de aproximações de longa e curta distância bem como a aproximação de curta distância e a aproximação de proximidade. O controle de atitude segue estratégias sob o mesmo prisma de movimento absoluto e relativo.


% \subsection{Lançamento: vínculos associados às missões de \gls{rvdb}}
% \label{sec:lancamento_vinculos}

% O início de uma missão de \gls{rvdb} ocorre na etapa de lançamento, conforme ilustrado na Figura~\ref{fig_etapas_rvdb}.
% Nessa fase, a espaçonave perseguidora deve ser injetada em uma órbita compatível com a do alvo, considerando a janela de lançamento planejada, de modo a evitar inserções que demandem correções adicionais de trajetória.

% A definição da janela de lançamento é vinculada à órbita conhecida do alvo.
% Os parâmetros de injeção são escolhidos em função dessa órbita, com o objetivo de garantir a inserção do veículo em uma trajetória coplanar à do alvo, reduzindo a necessidade de manobras posteriores para correção do plano orbital.

%%%%%%%%%%%%%%%%%%%%%%%%%%%%%%%%%%%%%%%%%%%%%%%%%%%%%%%%%%%%%%%%

% \section{Estratégia de controle de atitude}
% \label{sec:controle_atitude}


% O sistema de controle opera em dois modos distintos:
% (1) quando o manipulador robótico estiver travado (descrito pela dinâmica de corpo rígido) e
% (2) quando o \gls{mr} estiver em operação (descrito pela dinâmica acoplada da atitude e variáveis do movimento do \gls{mr}).
% Na fase de aproximação final, uma lei de controle atua nas juntas do manipulador para executar a manobra de captura.
% A estratégia inclui também os pontos de referência (\emph{waypoints}), que são pontos de espera para possíveis ajustes de posição e velocidades e cheques de segurança do sistema robótico.


%%%%%%%%%%%%%%%%%%%%%%%%%%%%%%%%%%%%%%%%%%%%%%%%%%%%%%%%%%%%%%%%
% \subsection{Controle PD na dinâmica relativa}
% \label{sec:controle_PD_HCW}

% A aproximação relativa no referencial \gls{lvlh}, definido no centro de massa da espaçonave alvo, é descrita pelo modelo linearizado \gls{hcw}.
% Neste referencial é descrito o movimento relativo entre a espaçonave perseguidora e a espaçonave alvo.
% A abordagem de controle do movimento orbital relativo é realizado pela lei de controle \gls{pd}.
% O objetivo é conduzir o movimento relativo a um órbita cuja deriva possa ser utilizada para reduzir um ângulo de fase que atenda a condição de aproximação final com o \gls{mr} em modo não operacional. 

%%%%%%%%%%%%%%%%%%%%%%%%%%%%%%%%%%%%%%%%%%%%%%%%%%%%%%%%%%%%%%%%%%%%%%%%%%%%%%%%%%

% \subsection{Controle de atitude durante o modo não operacional do \gls{mr}}
% \label{sec:controle_atitude_Euler}

% No modo não operacional, a \gls{etmr} é tratada como corpo rígido.
% A estratégia consiste em aplicar a técnica \gls{pd} para controle da atitude descrita em quaternion de erro entre a orientação da espaçonave perseguidora e uma orientação de referência associada ao alvo medida no sistema corpo.
% O controle é responsável por manter o perseguidor alinhado com o alvo durante a aproximação curta, enquanto o controle translacional define a trajetória relativa.

%%%%%%%%%%%%%%%%%%%%%%%%%%%%%%%%%%%%%%%%%%%%%%%%%%%%%%%%%%%%%%%%%%%%%%%%%%%%%%

% \subsection{Controle de atitude durante o modo operacional do \gls{mr}}
% \label{sec:controle_atitude_Lagrange}
% % 
% No modo operacional, o \gls{mr} passa a se movimentar para engatar o \gls{eef} na porta de acoplamento.
% A dinâmica do sistema deixa de ser equivalente à de um corpo rígido e passa a ser descrita pela formulação de Lagrange com coordenadas generalizadas que incluem a atitude da base e as juntas do \gls{mr}.
% Os movimentos articulares geram torques internos que aparecem como reações na base, alterando a sua atitude.

% A estratégia de controle utiliza o mesmo controlador que estabiliza a atitude quando o \gls{mr} está travado, e então passa a atuar também como compensador das reações provenientes dos movimentos das juntas.
% Assim, o controle de atitude deve ao mesmo tempo seguir a referência de atitude em relação ao alvo e anular os efeitos dos torques e forças de reação na plataforma robótica.

%%%%%%%%%%%%%%%%%%%%%%%%%%%%%

% \subsection{Síntese das estratégias de controle}


% Por fim, a Tabela~\ref{tab_locais_atuacao_controladores} apresenta a forma de implementação da lei de controle \gls{pd} neste trabalho, onde e como cada uma atua, de forma a cumprir a missão de \gls{rvdb}.

% \begin{table}[H]
% \centering
% \caption{Estratégia da lei de controle implementada.}
% \label{tab_locais_atuacao_controladores}
% \begin{tabular}{p{0.32\linewidth} p{0.62\linewidth}}
% \hline
% \textbf{Local de atuação} & \textbf{Como atua} \\
% \hline
% Controle de órbita &
% Reduz o erro e mantém a acurácia da posição relativa entre as espaçonaves e translada o perseguidor para próximo do alvo durante a aproximação curta. \\
% \hline
% Atitude relativa &
% Durante as aproximações curta distância e final, mantém a atitude do perseguidor sincronizada com a do alvo. \\
% \hline
% Juntas do \gls{mr} &
% Realiza a movimentação do \gls{mr} para que o efetuador final seja conduzido até a porta de acoplamento, de maneira suave e controlada. \\
% \hline
% \end{tabular}
% \end{table}